\chapter{Conclusion}
This thesis investigated how large language models can be used to support natural language-based comprehension and reprogrammability of robotic programs. The work contributed an architecture and implementation of an LLM-assisted system that combines retrieval-augmented generation, program/document indexing, and a user interface for exploration and modification of Universal Robots robotic programs.

The evaluation shows a mixed but informative outcome. For comprehension, the proposed system preserved near-perfect correctness (RQ1), indicating that retrieval-supported assistant responses and structured program views can help users identify key program properties even when working with low-level representations. However, comprehension time increased (RQ2), and confidence was lower than in the baseline condition. A significant part of the time increase is linked to assistant latency and interaction overhead rather than task complexity alone.

For modification and reuse, the findings are less favorable. The proposed system reduced modification correctness and increased completion time (RQ3 and RQ4), especially for tasks requiring structural edits such as node insertion. Participants were more successful on straightforward parameter updates than on transformations requiring deeper structural understanding of raw \texttt{.urpx} data. This indicates that the current workflow and representation are not yet a suitable foundation for reliable editing support in practical settings.

Together, the results answer the main research question as follows: LLMs can support natural language-based robotic program comprehension, but robust support for reprogrammability requires stronger grounding, better domain coverage, and a better editing abstraction than the one implemented in this thesis.

Beyond performance comparisons, the thesis contributes empirical insight into design pitfalls for robotics-oriented LLM systems. In this domain, response quality is strongly dependent on knowledge-base coverage and retrieval quality, and usability is highly sensitive to latency. The results also highlight that system-level workflow design matters as much as model capability: a chat assistant layered on top of a low-level program representation can help users reason about code, but it does not automatically make complex edits easier.

In conclusion, this thesis provides a practical baseline for continued work in LLM-assisted robotic programming. The findings suggest that future systems should combine improved retrieval and data quality with either higher-level editing abstractions or specialized model support for program transformations. With these improvements, the approach has clear potential to reduce barriers to robot program comprehension, modification, and reuse.
